\begin{recipe}{Basiskoekjes}
\kicker[Bakken,Zoet,Kinderen]{Bakken · zoet · kinderen}
\index[register]{Basiskoekjes}

\meta{45 minuten \dvd Ca.\ 25 stuks \dvd Oven 190\,°C}

\lettrine{E}{en} simpel koekjesdeeg dat zich goed leent om samen met een peuter
te maken. Het deeg is makkelijk te kneden en de vormpjes uitsteken is het leukste
gedeelte.

\blockrule{Ingrediënten}
\begin{ingredients}
  \ing{200 g}{patentbloem}
  \ing{100 g}{koude boter, in stukjes (50\%)}
  \ing{75 g}{kristalsuiker (37,5\%)}
  \ing{1}{ei}
  \ing{1 snufje}{zout}
\end{ingredients}

\blockrule{Bereiding}
\tip{Laat de koekjes goed afkoelen op het rooster voor je ze proeft. Uit de oven
zijn ze nog zacht, pas na het afkoelen worden ze knapperig.}
\begin{steps}
  \item Meng de bloem, suiker en het zout in een kom. Maak een kuiltje in het
        midden en breek het ei erin.
  \item Voeg de koude boter in stukjes toe en kneed alles door elkaar met de
        deeghaken van een mixer of met de hand, tot een egaal deeg.
  \item Wikkel het deeg in vershoudfolie en laat 30 minuten rusten in de koelkast.
        Zo wordt de boter weer koud, wat de koekjes tijdens het bakken minder
        laat uitzakken.
  \item Verwarm de oven voor op 190\,°C (boven- en onderwarmte). Bestrooi het aanrecht licht met bloem
        en rol het deeg uit tot 5--10 mm dikte.
  \item Steek vormpjes uit of snijd ze met een mes. Leg ze op een met bakpapier
        beklede bakplaat.
  \item Bak 12 minuten tot de koekjes goudbruin zijn. Laat volledig afkoelen op
        een rooster.
\end{steps}
\end{recipe}
